\begin{table}[tb]
\centering
\caption{Classifying features of roles, extracted
from~\autocite{steimann2000b,kuehn2014}.
\label{tab-features}}
\begin{tabular}[]{@{}cp{0.87\columnwidth}}
\toprule
\textbf{Nr.}
& 
\textbf{Feature of Roles \emph{(Meta level)}}
\tabularnewline
\midrule

1.
 & 
Roles have properties and behaviors \emph{(M1, M0)}
\tabularnewline

2.
 & 
Roles depend on relationships \emph{(M1)}
\tabularnewline

3.
 & 
Objects may play different roles simultaneously \emph{(M1, M0)}
\tabularnewline

4.
 & 
Objects may play the same role (type) several times \emph{(M0)}
\tabularnewline

5.
 & 
Objects may acquire and abandon roles dynamically \emph{(M0)}
\tabularnewline

6.
 & 
The sequence of role acquisition and removal may be restricted
\emph{(M1, M0)}
\tabularnewline

7.
 & 
Unrelated objects can play the same role \emph{(M1)}
\tabularnewline

8.
 & 
Roles can play roles \emph{(M1, M0)}
\tabularnewline

9.
 & 
Roles can be transferred between objects \emph{(M0)}
\tabularnewline

10.
 & 
The state of an object can be role-specific \emph{(M0)}
\tabularnewline

11.
 & 
Features of an object can be role-specific \emph{(M1)}
\tabularnewline

12.
 & 
Roles restrict access \emph{(M0)}
\tabularnewline

13.
 & 
Different roles may share structure and behavior \emph{(M1)}
\tabularnewline

14.
 & 
An object and its roles share identity \emph{(M0)}
\tabularnewline

15.
 & 
An object and its roles have different identities \emph{(M0)}
\tabularnewline

16.
 & 
Relationships between roles can be constrained \emph{(M1)}
\tabularnewline

17.
 & 
There may be constraints between relationships \emph{(M1)}
\tabularnewline

18.
 & 
Roles can be grouped and constrained together \emph{(M1)}
\tabularnewline

19.
 & 
Roles depend on compartments \emph{(M1, M0)}
\tabularnewline

20.
 & 
Compartments have properties and behaviors \emph{(M1, M0)}
\tabularnewline

21.
 & 
A role can be part of several compartments \emph{(M1, M0)}
\tabularnewline

22.
 & 
Compartments may play roles like objects \emph{(M1, M0)}
\tabularnewline

23.
 & 
Compartments may play roles which are part of themselves \emph{(M1,
M0)}
\tabularnewline

24.
 & 
Compartments can contain other compartments \emph{(M1, M0)}
\tabularnewline

25.
 & 
Different compartments may share structure and behavior
\emph{(M1)}
\tabularnewline

26.
 & 
Compartments have their own identity \emph{(M0)}
\tabularnewline
\bottomrule
\end{tabular}
\end{table}
